% source: https://tex.stackexchange.com/a/764195 (question 762384, score 8, block 1)
% Source - https://tex.stackexchange.com/a/762133
% Posted by PHL, modified by community. See post 'Timeline' for change history
% Retrieved 2026-04-22, License - CC BY-SA 4.0

% !TEX TS-program = LuaLaTeX
\documentclass[tikz]{standalone}
\usepackage[svgnames]{xcolor}
\usepackage{luadraw}

\makeatletter
\def\luadrawLiteral#1{\pgfsys@invoke{#1}}% writes in pf stream
\makeatother

\begin{luacode*}
local ld = luadraw
local cpx = ld.cpx

function ld.graph:bpCoord(z) -- z=coordinates, conversion cm -> bp
    local a, b = self:coord(z.re,z.im)
    local c = 28.3464567
    return  ld.strReal(a*c,2).." "..ld.strReal(b*c,2) -- two decimals
end

function ld.graph:Dliteralpolyline(...) -- writes polylines in the enter image description herepdf stream
-- Dliteralpolyline( polyline1, args1, polyline2, args2, ...)
-- args = {fill= rgb table, color = rgb table, close = false}, (fill is solid)
-- fill="none" => no fill, color="none" => no line 
    local fill, draw, ends = {}, {}
    local close = false
    local literalstr = {}
    local strcolor, strfill, strdraw 
    
    local function apolyline(polyline,args)
        if (type(polyline[1]) == "number") or cpx.isComplex(polyline[1]) then
            polyline = {polyline}
        end
        fill = args.fill or fill
        draw = args.draw or draw
        if args.close ~= nil then close = args.close end
        if (draw == "none") and (fill == "none") then return -- no fill, no stroke
        elseif (draw == "none") then ends = "f" -- only fill
        elseif (fill == "none") then ends = "s" -- only stroke
        else ends = "b" -- fill and stroke
        end 
        if (type(fill) == "table") and (#fill == 3) then
            strcolor = ld.strReal(fill[1],3).." "..ld.strReal(fill[2],3).." "..ld.strReal(fill[3],3)
            if strcolor ~= strfill then
                table.insert(literalstr,strcolor.." rg" ) -- fill color
                strfill = strcolor
            end
        end
        if (type(draw) == "table") and (#draw == 3) then
            strcolor = ld.strReal(draw[1],3).." "..ld.strReal(draw[2],3).." "..ld.strReal(draw[3],3)
            if strcolor ~= strdraw then
                table.insert(literalstr,strcolor.." RG" ) -- line color
                strdraw = strcolor
            end
        end
        local aux = {}
        local L = self:Mtransform( polyline )
        local last
        for _,cp in ipairs(L) do
            for k,z in ipairs(cp) do 
                z = cpx.toComplex(z)
                if k == 1 then table.insert(literalstr, self:bpCoord(z).." m") 
                else
                    local bpcoord = self:bpCoord(z)
                    if  last ~= bpcoord then table.insert(literalstr, bpcoord.." l") end
                    last = bpcoord
                end
            end
            if close then table.insert(literalstr, "h") end
        end
        table.insert(literalstr, ends)
    end
    
    local n = select("#", ...)  -- Nombre total d'arguments
    for i = 1, n-1, 2 do   -- Pas de 3 (1,4,7...)
        local polyline, args = select(i, ...)  -- Récupère les 3 args
        if polyline ~= nil then apolyline(polyline,args) end
    end
    self:Writeln("\\luadrawLiteral{"..table.concat(literalstr,"\n").."}")
end
\end{luacode*}

\directlua{t1 = os.clock()}%

\begin{document}
\begin{luadraw}{name=penrose, exec=true}
local ld = luadraw
local cpx = ld.cpx
local Z, Zp = cpx.Z, cpx.Zp
local iterations = 8

 -- golden ratio, its inverse and it's inverse squared
local phi, invphi, invphisq = (1+math.sqrt(5))/2, 2/(1+math.sqrt(5)), (2/(1+math.sqrt(5)))^2

local r = {} -- to store the tenth roots of unity
for k = 0, 9 do
    r[k] = Zp(1, k*math.pi/5)
end

local S = {} -- store the Robinson triangles
-- we start with the wheel
for k = 1, 8 do
    if (k % 2 == 0) then
        S[k] = {4, 0, r[k], r[k+1]}
    else
        S[k] = {3, 0, r[k+1], r[k]}
    end
end
S[9] = {3,0, 1, r[9]}
S[10] = {4,0, 1, r[1]}

-- {n, v1, v2, v3}
-- 1 = counter clockwise gnomon
-- 2 = clockwise gnomon
-- 3 = counter clockwise golden triangle
-- 4 = clockwise golden triangle
-- v1 = the only one with a different angles, then follows the given direction
-- replacement procedure for Robinson triangles
local function selfsim(R)
    local t, a, b, c = table.unpack(R)
    local d = b*invphi + c*invphisq
    local e = c*invphi + a*invphisq
    local f = a*invphi + b*invphisq
    if t == 1 then
        return {{1, d, a, b}, {4, c, d, a}}
    elseif t == 2 then
        return {{2, d, a, b}, {3, c, d, a}}
    elseif t == 3 then
        return {{4, b, f, e}, {3, b, c, e}, {1, f, e, a}}
    else
        return {{3, b, f, e}, {4, b, c, e}, {2, f, e, a}}
    end
end

-- iteratively construct the full list of triangles
-- construct the full list of tiles
for k = 1, iterations do -- iterate
    local newS = {} -- temp variable
    for _, s in ipairs(S) do -- for each j element of S
        for t, v in ipairs(selfsim(s)) do -- do selfsim(j), obtains 2 or 3 triples, for each thus triple
            table.insert(newS, v) -- add it to the new list
        end
    end
    S = newS -- S is now the new list
end

-- start the graph
local x1,x2,y1,y2 = ld.getbounds(ld.map(function(v) return {v[2],v[3],v[4]} end, S))
local sc = math.sqrt(2)
local g = ld.graph:new{ window={x1,x2,y1,y2}, size={10,10},margin={0,0,0,0}}

-- draw the polygons
local colors = {ld.LightBlue, ld.Orange, ld.ForestGreen, ld.Indigo}
local opacity = 1
local triangles = false
g:Lineoptions("solid","white",1)
g:Scale(sc)
local pdfliteral = {}
g:Filloptions(nil,nil, opacity) --only once to adjust the opacity
for j = 1, #S do
    local t, a, b, c = table.unpack(S[j])
    --g:Filloptions("full",colors[t], opacity) -- takes far too long
    if triangles then
        table.insert(pdfliteral, {a,b,c})
        table.insert(pdfliteral, {close=true, fill=colors[t]})
    else
        if t == 1 then -- t=1 <=> anti-clockwise gnomon
            table.insert(pdfliteral, {a, b, c, a + r[7] * (c - a)})
            table.insert(pdfliteral, {close=true, fill=colors[t]})
        elseif t == 4 then -- t=3 <=> anti-clockwise golden triangle
            table.insert(pdfliteral, {a, b, c, c + invphi*r[3]*(c - a)})
            table.insert(pdfliteral, {close=true, fill=colors[t]})
        end
    end
end
g:Dliteralpolyline( table.unpack(pdfliteral) )
g:Show()
\end{luadraw}

\directlua{t2 = os.clock(); print("\\ntime=",t2-t1)}%
\end{document}
